\chapter*{Q}
\label{chap:q}

\term{Q-analog}
A $q$-analog introduces a parameter $q$ into a formula so that the classical expression is recovered as $q\to1$. The $q$-binomial coefficient is a standard example. \par\noindent\textbf{Applications:} $q$-analogs occur in combinatorics, number theory, statistical mechanics, and quantum algebra.

\term{Q-matrix}
For a continuous-time Markov chain, a $Q$-matrix has nonnegative off-diagonal transition rates and row sums zero, so $q_{ii}=-\sum_{j\ne i}q_{ij}$. \par\noindent\textbf{Applications:} $Q$-matrices describe queues, reliability systems, population models, and chemical reaction networks.

\term{Q-series}
A $q$-series is a power series in $q$, often with coefficients encoding arithmetic or combinatorial data. The generating function for integer partitions is a famous example. \par\noindent\textbf{Applications:} $q$-series support partition theory, modular forms, combinatorics, and statistical mechanics.

\term{Quadratic Equation}
A quadratic equation has the form $ax^2+bx+c=0$ with $a\ne0$. Its roots are $x=(-b\pm\sqrt{b^2-4ac})/(2a)$, and the discriminant determines whether the roots are real or complex. \par\noindent\textbf{Applications:} quadratic equations model trajectories, optimization, geometry, and basic physical laws.

\term{Quadratic Form}
A quadratic form is a homogeneous degree-two expression $Q(v)=v^{\mathsf T}Av$, with $A$ taken symmetric. Its definiteness describes curvature and its level sets include ellipses, hyperbolas, and higher-dimensional analogues. \par\noindent\textbf{Applications:} quadratic forms appear in optimization, energy, statistics, geometry, and number theory.

\term{Quadratic Reciprocity}
For distinct odd primes $p$ and $q$, quadratic reciprocity states that $\left(\frac{p}{q}\right)\left(\frac{q}{p}\right)=(-1)^{\frac{p-1}{2}\frac{q-1}{2}}$. It relates whether $p$ is a square modulo $q$ to whether $q$ is a square modulo $p$. \par\noindent\textbf{Applications:} it makes quadratic-residue calculations efficient and supports number theory and cryptography.

\term{Quadratic Residue}
An integer $a$ is a quadratic residue modulo a prime $p$ if $x^2\equiv a\pmod p$ has a solution. Otherwise it is a nonresidue. \par\noindent\textbf{Applications:} quadratic residues support modular arithmetic, primality tests, cryptography, and number-theoretic reciprocity.

\term{Quadrature}
Quadrature means evaluating a definite integral, often by a weighted numerical sum. Gaussian quadrature chooses nodes and weights to integrate polynomials of high degree exactly. \par\noindent\textbf{Applications:} quadrature is used in simulation, probability, finite elements, and engineering computation.

\term{Quadrature Domain}
A quadrature domain is a bounded complex domain for which integrals of suitable holomorphic functions equal a finite combination of function values and derivatives at specified points. \par\noindent\textbf{Applications:} quadrature domains connect complex analysis with exact integration and potential-flow models.

\term{Quadrature Rule}
A quadrature rule approximates an integral by a weighted sum of function values, such as the trapezoidal, Simpson, or Gaussian rule. Accuracy depends on smoothness and node placement. \par\noindent\textbf{Applications:} quadrature rules evaluate integrals when elementary antiderivatives are unavailable.

\term{Qualitative}
Qualitative analysis describes properties such as stability, boundedness, or long-term behavior without finding an exact formula. \par\noindent\textbf{Applications:} it provides useful conclusions for dynamical systems, differential equations, and models that cannot be solved explicitly.

\term{Qualitative Analysis}
Qualitative analysis studies features such as stability, equilibria, limit cycles, and long-term behavior without finding an explicit formula. \par\noindent\textbf{Applications:} it analyzes ODEs, dynamical systems, population models, and physical systems that are not exactly solvable.

\term{Quantifier}
A quantifier specifies how widely a statement applies. The universal quantifier $\forall$ means “for every,” and the existential quantifier $\exists$ means “for at least one.” \par\noindent\textbf{Applications:} quantifiers express mathematical definitions, theorems, specifications, and computer-verification statements.

\term{Quantile}
A quantile is a cutoff dividing a distribution so that a specified fraction of observations lies below it. The $0.5$ quantile is the median, and the $0.25$ and $0.75$ quantiles define the quartiles. \par\noindent\textbf{Applications:} quantiles summarize skewed data, define thresholds, and support robust statistics.

\term{Quantile Function}
The quantile function is $Q(p)=\inf\{x:F(x)\geq p\}$, where $F$ is a cumulative distribution function. It returns a value below which a fraction $p$ of the distribution lies. \par\noindent\textbf{Applications:} quantile functions support simulation, risk thresholds, confidence intervals, and robust summaries.

\term{Quantitative}
Quantitative information is expressed through numbers, measurements, or numerical relationships. \par\noindent\textbf{Applications:} quantitative descriptions support modeling, comparison, prediction, experiments, and data-driven decision-making.

\term{Quantitative Analysis}
Quantitative analysis uses numerical data, mathematical models, and statistical methods to describe patterns and make predictions. \par\noindent\textbf{Applications:} it supports finance, engineering, epidemiology, experimentation, and data-driven policy.

\term{Quantity}
A quantity is a measurable or countable amount, such as length, mass, probability, or number. Quantities may be represented by numbers with units or by abstract algebraic elements. \par\noindent\textbf{Applications:} quantities provide the variables and parameters used in scientific models and measurements.

\term{Quantum State}
A quantum state is a unit vector, or physically a ray, in a Hilbert space. Measurement probabilities are determined by squared inner products, such as $|\langle\psi\mid\phi\rangle|^2$. \par\noindent\textbf{Applications:} quantum states describe qubits, atoms, molecules, quantum computation, and quantum communication.

\term{Quarter}
A quarter is one fourth of a whole. In the plane, a quadrant is one of the four regions determined by the coordinate axes. \par\noindent\textbf{Applications:} quarters and quadrants organize fractions, angles, coordinate geometry, and seasonal or financial data.

\term{Quarter-turn Symmetry}
Quarter-turn symmetry means that a $90$-degree rotation leaves an object unchanged. Repeating the rotation gives the cyclic group of order four. \par\noindent\textbf{Applications:} this symmetry appears in tilings, logos, crystal patterns, image processing, and geometric design.

\term{Quasicoherent Sheaf}
A quasicoherent sheaf is a sheaf of modules that locally corresponds to a module over the coordinate ring of an affine piece. It is the natural sheaf-theoretic form of algebraic data on a scheme. \par\noindent\textbf{Applications:} quasicoherent sheaves describe functions, ideals, vector bundles, and cohomology in algebraic geometry.

\term{Quasiconcave Function}
A function is quasiconcave when every superlevel set $\{x:f(x)\geq\alpha\}$ is convex. It is useful for functions whose high-value regions have no gaps. \par\noindent\textbf{Applications:} quasiconcavity appears in economics, utility theory, location problems, and optimization.

\term{Quasiconformal Map}
A quasiconformal map is a homeomorphism that may distort angles but has uniformly bounded eccentricity. Infinitesimal circles become ellipses whose axis ratio is controlled. \par\noindent\textbf{Applications:} quasiconformal maps are used in complex analysis, geometric topology, image registration, and shape deformation.

\term{Quasi-convex Function}
A function is quasiconvex when every sublevel set $\{x:f(x)\leq\alpha\}$ is convex. Every convex function is quasiconvex, but not conversely. \par\noindent\textbf{Applications:} quasiconvexity supports optimization when objectives are not strictly convex.

\term{Quasigroup}
A quasigroup is a set with a binary operation for which both $a*x=b$ and $y*a=b$ have unique solutions. It need not be associative or have an identity. \par\noindent\textbf{Applications:} quasigroups appear in algebraic combinatorics, Latin squares, coding, and cryptographic constructions.

\term{Quasi-isomorphism}
A quasi-isomorphism is a map of chain complexes that induces isomorphisms on all homology groups. Derived categories invert these maps, treating complexes with the same homological information as equivalent. \par\noindent\textbf{Applications:} quasi-isomorphisms simplify homological algebra, derived geometry, and algebraic topology.

\term{Quasilinear}
A quasilinear differential equation is linear in its highest-order derivatives but may be nonlinear in lower-order terms or coefficients. \par\noindent\textbf{Applications:} quasilinear equations model fluid flow, waves, transport, and nonlinear continuum mechanics.

\term{Quasi-metric Space}
A quasi-metric space has a distance-like function satisfying nonnegativity, identity, and the triangle inequality but not necessarily symmetry. Thus traveling from $x$ to $y$ may have a different cost from traveling back. \par\noindent\textbf{Applications:} quasi-metrics model directed networks, transportation costs, computer science, and asymmetric geometry.

\term{Quasi-Newton Method}
A quasi-Newton method updates an approximation to the Hessian or its inverse using changes in gradients, rather than computing second derivatives directly. BFGS is a standard example. \par\noindent\textbf{Applications:} quasi-Newton methods solve smooth optimization problems in statistics, machine learning, and engineering.

\term{Quasi-periodic Function}
A quasiperiodic function combines oscillations with incommensurate frequencies. It does not repeat exactly, but its values can return arbitrarily close to previous states. \par\noindent\textbf{Applications:} quasiperiodicity models celestial motion, oscillators, solid-state systems, and dynamical systems.

\term{Quasi-random Sequence}
A quasirandom sequence is a low-discrepancy sequence designed to cover a region uniformly rather than randomly. Sobol and Halton sequences are common examples. \par\noindent\textbf{Applications:} quasirandom sequences improve numerical integration, uncertainty quantification, finance, and computer graphics.

\term{Quaternion Algebra}
A quaternion algebra is a four-dimensional central simple algebra over its center that generalizes the usual quaternions. Its multiplication may be described by generators and relations, and it need not be a division algebra over every field. \par\noindent\textbf{Applications:} quaternion algebras are used in number theory, quadratic forms, arithmetic geometry, and the study of rotations.

\term{Quaternion Norm}
A quaternion norm is defined for $q=a+bi+cj+dk$ by $|q|=\sqrt{a^2+b^2+c^2+d^2}$. It is multiplicative: $|q_1q_2|=|q_1||q_2|$. \par\noindent\textbf{Applications:} the norm supports inversion, numerical stability, and rotation computations.

\term{Quaternions}
A quaternion is an element of the noncommutative division algebra with basis $1,i,j,k$ and relations $i^2=j^2=k^2=ijk=-1$. Unit quaternions represent spatial rotations without the singularities of some angle-coordinate systems. \par\noindent\textbf{Applications:} quaternions are used in computer graphics, robotics, aerospace, and physics.

\term{Quine}
A quine is a program that prints its own source code without reading it as external input. It is constructed using quotation and self-reference. \par\noindent\textbf{Applications:} quines illustrate recursion, fixed-point theorems, programming-language semantics, and computability.

\term{Quintic}
A quintic is a polynomial of degree five. Some quintics are solvable by radicals, but the general quintic is not. \par\noindent\textbf{Applications:} quintics illustrate polynomial classification and the role of Galois groups.

\term{Quintic Equation}
A quintic equation is a polynomial equation of degree five. The Abel--Ruffini theorem states that no formula using radicals solves every quintic, although particular quintics can be solvable. \par\noindent\textbf{Applications:} quintics illustrate the limits of symbolic algebra and the role of Galois theory in polynomial equations.

\term{Quotient Group}
If $N$ is normal in $G$, the quotient group $G/N$ consists of cosets with $(aN)(bN)=(ab)N$. It treats elements differing by $N$ as equivalent. \par\noindent\textbf{Applications:} quotient groups describe reduced symmetries, congruences, and the structure of homomorphisms.

\term{Quotient Manifold}
A quotient manifold is formed by identifying points of a manifold under a free, properly discontinuous group action. The quotient inherits a smooth structure, as in a torus obtained from a plane by lattice translations. \par\noindent\textbf{Applications:} quotient manifolds model periodic geometry, covering spaces, and symmetry reductions.

\term{Quotient Map}
A quotient map is a surjective map $f:X\to Y$ such that $U\subseteq Y$ is open exactly when $f^{-1}(U)$ is open in $X$. It creates the quotient topology and is continuous, but need not be open. \par\noindent\textbf{Applications:} quotient maps construct circles, projective spaces, and glued topological spaces.

\term{Quotient Ring}
For a two-sided ideal $I$ of a ring $R$, the quotient ring $R/I$ consists of cosets $r+I$ with induced addition and multiplication. For example, $\mathbb C\cong\mathbb R[x]/(x^2+1)$. \par\noindent\textbf{Applications:} quotient rings encode congruences, construct fields, and define algebraic varieties.

\term{Quotient Space}
A quotient space $X/{\sim}$ is formed by identifying points related by an equivalence relation. The quotient map sends each point to its equivalence class. \par\noindent\textbf{Applications:} quotient spaces construct tori, projective spaces, and spaces obtained by gluing boundaries.

\term{Quotient Topology}
The quotient topology on $Y=X/{\sim}$ is the finest topology for which the quotient map $q:X\to Y$ is continuous; equivalently, $U\subseteq Y$ is open when $q^{-1}(U)$ is open. \par\noindent\textbf{Applications:} it formalizes gluing and identification in topology and geometry.
